\subsection*{Supporting analyses (original design)}
The following analyses used the single-split design and the mean-viability GDSC2
summary of the earlier version of this study, in which concordance was computed on
all shared cell lines; they are reported as supporting evidence, and the primary
analyses in the main text supersede them where the two overlap. (i) Determinants: each drug
was described by the dynamic range of its sensitivity distribution in each assay
and by the lineage diversity of its sensitive cell lines (Shannon entropy of DepMap
OncotreeLineage in the CTRP-sensitive tertile), related to concordance by Spearman
correlation and a standardized multivariable regression. (ii) Reverse direction:
models trained on GDSC2 labels and tested on CTRP. (iii) Continuous formulation: a
random-forest regressor trained on the untransformed CTRP value, with transfer
measured as the Spearman correlation between predictions and GDSC2 sensitivity of
held-out cell lines. (iv) Sensitivity to median, tertile and quartile labels, 50 to
500 selected features, and univariate, mutual-information and L1-penalized feature
selection. (v) Jaccard overlap of the top-200 genes selected on CTRP and on GDSC2
labels, mechanism-of-action class, and the rank of canonical target genes.

\subsection*{Exploratory extension to patient tumors}
Each drug's model was trained on all CCLE cell lines carrying CTRP tertile-extreme
labels and applied to tumors from The Cancer Genome Atlas (TCGA) whose patients
received the same drug. Patient expression was the Genomic Data Commons log-TPM
protein-coding matrix, restricted to the 19{,}125 genes shared with CCLE through
Ensembl identifiers; distribution shift was reduced by standardizing each gene
within each domain separately, which uses the distribution of each drug's patient
set. Patient outcome was the clinical drug response recorded in TCGA, binarized as
responder (complete or partial response) versus non-responder (stable or
progressive disease) (Ding \emph{et al.}, \emph{Bioinformatics} 2016;32:2891--2895); when a patient had several records
for a drug, the first was used. Drugs with at least 25 evaluable patients and five
per response class were retained. Treatment context, combination therapy and line
of therapy were not modeled, so this analysis is exploratory.

\subsection*{Leave-one-assay-out labels}
For each ordered pair of assays, the training assay's candidate cell lines were its
tertile extremes over all its cell lines; after the 70/30 split, training thresholds
were re-estimated without the test cell lines. Labels in the receiving assay were
defined by tertiles of that assay's responses among the cell lines not used for
training. The held-out assay's labels therefore never enter the score, but its
thresholds are set on the evaluated population rather than fixed in advance.

\subsection*{Single-reference-set analysis}
For each of the ten saved splits, each drug received one random subset of $k$ = 20,
50 or 100 cell lines from its reference set (cell lines with CTRP and fitted GDSC2
responses, excluding that split's external test cell lines). Concordance on that
subset was correlated across drugs with the same split's external random-forest
AUROC; the screening ROC-AUC (external AUROC $\geq 0.70$) and the agreement of the
$\rho \geq 0.5$ call with the full-reference call were also recorded. Each size was
drawn 200 times per split (2{,}000 replicates; seed 42; no model retrained).
Results are in Table~S9 and main Figure~3B.

\subsection*{Use of AI tools}
Claude (Anthropic), used through the Claude Code command-line tool with Claude
models available between July and October 2026, was used to write and run the
analysis code, generate the figures and supplementary tables, draft and revise the
manuscript text, check reported numbers against the result files, and edit
language. The tool worked in the author's computing environment on the public input
data listed in the main text; no patient-identifiable or non-public data were
provided to it. The following safeguards were applied. Every headline number in the
main text is checked by a script (\texttt{verify\_numbers\_rev.py}, 67 checks)
that recomputes it from the saved result files and requires the quoted phrase to
appear in the manuscript. Scripts re-run from the public code repository, including
those for the single-reference-set, drug-block permutation and valid-split analyses,
reproduced the corresponding result files used in the paper byte for byte, and the
supplement built from the repository is identical to this document. Analyses,
figures and text were reviewed by the author, who takes full responsibility for the
content.

